% !TEX root = ../main.tex
\chapter{Uso dell'Intelligenza Artificiale}
\label{cap:processo-ia}

Durante la realizzazione del progetto e della documentazione sono stati utilizzati strumenti di intelligenza artificiale per ottimizzare i tempi, supportare la configurazione dell'infrastruttura e la risoluzione di problemi. L'intero processo ha previsto una costante supervisione umana: ogni porzione di codice o testo è stata verificata, adattata e validata manualmente prima dell'integrazione.

\section{Supporto allo sviluppo software}

Nella fase di implementazione, l'uso dell'IA ha riguardato prevalentemente attività ripetitive e di diagnostica:
\begin{itemize}
    \item \textbf{Generazione di codice e test:} stesura di porzioni di codice ripetitivo e generazione della struttura base dei test, successivamente completati e ricontrollati manualmente.
    \item \textbf{Debugging:} ricerca, individuazione e risoluzione di bug durante la fase di sviluppo.
    \item \textbf{Supporto tecnologico:} ricerca, spiegazione e affiancamento nella configurazione di librerie e strumenti specifici integrati nel progetto, in particolare HikariCP, Flyway, JUnit Jupiter, Mockito e JaCoCo.
\end{itemize}

\section{Supporto alla redazione della relazione}

L'apporto dell'IA per la stesura di questo documento si è concentrato su aspetti formali e strutturali:
\begin{itemize}
    \item \textbf{Infrastruttura LaTeX:} creazione della struttura di base del documento, configurazione dei comandi e risoluzione degli errori di compilazione.
    \item \textbf{Revisione linguistica:} correzione ortografica, sistemazione di imprecisioni grammaticali e uniformazione dello stile espositivo.
    \item \textbf{Bibliografia e tecnologie:} supporto nella ricerca e stesura dei riferimenti tecnologici (sia per gli strumenti proposti dall'IA, sia per quelli scelti autonomamente).
\end{itemize}

\section{Validazione e controllo}

L'apporto dell'IA è stato impiegato esclusivamente come acceleratore del processo. Tutto il materiale generato (codice sorgente, configurazioni, test o testo) è stato sistematicamente esaminato e ricontrollato dal team. Il gruppo ha mantenuto la piena responsabilità sulle decisioni architetturali e sull'implementazione delle regole di dominio.